\chapter{Univers, symbologie et opérations sur titres}\label{ch:rs:universe-symbology-and-corporate-actions}

Avant l'ouverture du marché le 9 juin 2022, les actions de classe A de la société auparavant appelée Facebook commencent à se négocier sur
le Nasdaq sous un nouveau ticker, META~; le ticker FB, utilisé depuis la cotation en 2012, est retiré. Un tableau de recherche
indexé sur les tickers perd dix ans d'historique à cette date ou, si l'ancien ticker est attribué plus tard à quelqu'un d'autre, raccorde deux
sociétés en une seule série. Les tickers ne sont pas la seule chose qui bouge~: les sociétés s'introduisent en bourse et en sortent, divisent leurs actions, fusionnent
et procèdent à des scissions, et l'ensemble des titres qu'une stratégie aurait pu négocier change chaque mois. Ce chapitre construit les objets qui
permettent de garder le contrôle de tout cela~: un univers défini dans le temps, un \omterm{def:rs:universe-symbology-and-corporate-actions:master}{référentiel des titres} indexé par des \omterm{def:rs:universe-symbology-and-corporate-actions:permanent}{identifiants permanents}, des ajustements
qui ne laissent pas fuir l'information future, et les \omterm{def:rs:universe-symbology-and-corporate-actions:delisting}{rendements de radiation}. Ses données sont un marché simulé d'environ 1\,500 titres cotés sur dix ans,
dans lequel les tickers changent et sont réattribués, les actions se divisent, et des sociétés font faillite ou sont rachetées.

\section{Un univers investissable dans le temps}

\begin{definition}[Univers investissable, appartenance à l'univers, filtre de liquidité]\label{def:rs:universe-symbology-and-corporate-actions:universe}
Un \emph{univers investissable}\index{univers investissable} est l'ensemble des titres qu'une stratégie peut détenir à une date donnée, défini par
des règles évaluées sur les données connues à cette date. L'\emph{appartenance à l'univers}\index{appartenance à l'univers} est le tableau qui en résulte, constitué de
couples (date, titre). Un \emph{filtre de liquidité}\index{filtre de liquidité} est une règle d'appartenance portant sur la négociabilité~: une valeur minimale de la
valeur moyenne quotidienne échangée, un cours minimal, une capitalisation minimale du flottant, ou un classement selon l'une de ces grandeurs.
\end{definition}

Un univers est une décision de recherche aux conséquences importantes. Il fixe les titres sur lesquels les prédicteurs sont mesurés (chapitre~6), ceux
que l'optimiseur peut acheter (chapitre~25) et le capital que la stratégie peut mobiliser (chapitre~28). Deux propriétés le rendent honnête~:
ses règles n'utilisent que des données connues à chaque date, et il ne change pas davantage que ne change la liquidité sous-jacente. La seconde
n'est pas automatique. Une règle qui retient les 500 titres à la plus forte valeur échangée chaque mois remplace chaque titre dont le rang
oscille autour de 500, alors même que rien n'a changé pour ce titre~; chaque remplacement constitue une paire de transactions que la stratégie
doit ensuite payer.

\begin{method}[Une règle d'appartenance à tampon]\label{met:rs:universe-symbology-and-corporate-actions:buffer}
À chaque date de rééquilibrage, classer les candidats selon la mesure de liquidité connue ce jour-là. Un membre conserve sa place tant que son rang
reste dans un rang de sortie plus large que la taille cible (par exemple 600 pour un univers de 500)~; les places libres vont aux non-membres les mieux
classés. Les fournisseurs d'indices utilisent le même dispositif (règles de tampon, livre~1, chapitre~15).
\end{method}

Dans le marché simulé du chapitre, un univers constitué des 500 titres à la plus forte valeur échangée remplace 11.5~\% de ses membres chaque mois
sans tampon, et 3.1~\% avec un rang de sortie de 600 (\cref{fig:rs:universe-symbology-and-corporate-actions:churn}). L'univers doté
d'un tampon n'est pas moins liquide sur un plan qui compte~: tous ses membres figurent parmi les 600 premiers.

\begin{omfigure}
\begin{tikzpicture}
\begin{axis}[width=11cm,height=4.8cm,xlabel={mois},ylabel={membres remplacés},tick label style={font=\small},
  label style={font=\small},xmin=0,xmax=120,ymin=0,ymax=100,grid=major,grid style={gray!25},legend columns=2,
  legend style={font=\footnotesize,at={(0.5,-0.32)},anchor=north,draw=none,fill=none,
  /tikz/every even column/.append style={column sep=8pt}},table/col sep=comma]
\addplot[omThm,thick] table[x=month,y=plain]{figdata/research/04-universe-symbology-and-corporate-actions/churn.csv};
\addplot[omDef,very thick] table[x=month,y=buffer]{figdata/research/04-universe-symbology-and-corporate-actions/churn.csv};
\legend{500 premiers chaque mois,500 premiers avec un rang de sortie de 600}
\end{axis}
\end{tikzpicture}

\omcaption{Membres d'un univers de liquidité de 500 titres remplacés à chaque rééquilibrage mensuel, dans le marché simulé du chapitre~:
11.5~\% par mois en moyenne sans tampon, 3.1~\% avec un tampon. Données~: module du chapitre, graine fixée.}\label{fig:rs:universe-symbology-and-corporate-actions:churn}
\end{omfigure}

\section{Des identifiants qui changent}

Le livre~1, chapitre~28, a posé la règle~: un ticker n'est pas un identifiant. Les tickers sont modifiés lorsqu'une société
change de nom, libérés lorsqu'elle est radiée, puis attribués plus tard à une nouvelle société. L'exemple n'est pas rare. Après que General Motors Corporation
a demandé la protection du chapitre~11 le premier juin 2009, ses actions se négocient de gré à gré (sous le symbole MTLQQ à la fin de cette année-là) et le ticker GM
était libre~; la nouvelle General Motors Company, qui avait racheté l'essentiel des actifs de l'ancienne société, a été cotée à la New York Stock
Exchange sous le ticker GM en novembre 2010. Un tableau indexé sur GM réunit donc deux sociétés différentes, séparées par une faillite.

Les identifiants nationaux sont plus stables sans être permanents. Un CUSIP est un code de 9 caractères pour les titres américains et canadiens, et un
ISIN (ISO 6166) est un code de 12 caractères composé d'un préfixe pays de deux lettres, d'un identifiant national et d'une clé de contrôle~; les deux
désignent une émission, et non une société, et relèvent des symboles externes qu'une entreprise doit faire correspondre plutôt que de les tenir pour fiables. Le Financial Instrument Global Identifier
(FIGI) est conçu pour ne pas changer~: une fois attribué, il n'est jamais modifié, et lorsque l'instrument cesse d'exister, le FIGI est
retiré et jamais réutilisé.

\begin{definition}[Identifiant permanent, correspondance des identifiants, changement de ticker]\label{def:rs:universe-symbology-and-corporate-actions:permanent}
Un \emph{identifiant permanent}\index{identifiant permanent} est un identifiant attribué à un titre lorsqu'il entre dans les données de
l'entreprise, et qui n'est jamais modifié ni réattribué, quoi qu'il arrive à ses symboles externes. Une \emph{correspondance des identifiants}\index{correspondance des identifiants}
est la table qui relie chaque symbole externe (ticker, CUSIP, ISIN, code du fournisseur) à un identifiant permanent sur un intervalle de dates.
Un \emph{changement de ticker}\index{changement de ticker} met fin à une ligne de cette correspondance et en ouvre une autre pour le même identifiant permanent.
\end{definition}

\begin{omfigure}
\begin{tikzpicture}[font=\footnotesize,x=1cm,y=0.9cm]
\draw[->,>=Stealth] (0,0) -- (12.6,0) node[right] {temps};
\foreach \x/\lab in {0/2008,2/2009,4/2010,6/2011,8/2012,10/2013,12/2014} \draw (\x,0.08) -- (\x,-0.08) node[below] {\lab};
\draw[omDef,line width=5pt] (0,3.0) -- (2.8,3.0);
\node[omDef,anchor=south west] at (0,3.25) {identifiant permanent 1044~: ticker GM};
\draw[omDef!50,line width=5pt] (2.8,3.0) -- (4.2,3.0);
\node[omDef!70,anchor=west] at (4.35,3.0) {1044~: de gré à gré};
\draw[omThm,line width=5pt] (5.8,1.6) -- (12.4,1.6);
\node[omThm,anchor=south west] at (5.8,1.85) {identifiant permanent 2877~: ticker GM};
\draw[black!60,dashed] (0,0.75) -- (2.8,0.75);
\draw[black!60,dashed] (5.8,0.75) -- (12.4,0.75);
\node[black!70,anchor=south west] at (6.9,0.8) {une table indexée sur GM~: deux sociétés, une seule série};
\draw[->,>=Stealth,black!60] (2.8,0.75) to[out=20,in=160] (5.8,0.75);
\end{tikzpicture}

\omcaption{Un ticker réattribué. La cotation d'\omterm{def:rs:universe-symbology-and-corporate-actions:permanent}{identifiant permanent} 1044 se négocie sous GM jusqu'à mi-2009, puis de gré à gré sous
un autre symbole~; à partir de fin 2010, une nouvelle cotation, 2877, se négocie sous GM. Un \omterm{def:rs:universe-symbology-and-corporate-actions:master}{référentiel des titres} conserve deux identifiants et deux intervalles de dates
~; une table indexée sur le ticker les réunit en une seule série avec une rupture. Les identifiants sont fictifs.}\label{fig:rs:universe-symbology-and-corporate-actions:reuse}
\end{omfigure}

\section{Le référentiel des titres}

\begin{definition}[Référentiel des titres]\label{def:rs:universe-symbology-and-corporate-actions:master}
Un \emph{référentiel des titres}\index{référentiel des titres} (security master) est la table de référence des titres de l'entreprise~: pour chaque \omterm{def:rs:universe-symbology-and-corporate-actions:permanent}{identifiant permanent}, son émetteur,
ses dates de cotation et de radiation, ses identifiants externes au fil du temps, ses opérations sur titres et sa raison de radiation ainsi que son
\omterm{def:rs:universe-symbology-and-corporate-actions:delisting}{rendement de radiation}~; toute autre table désigne les titres par leur seul \omterm{def:rs:universe-symbology-and-corporate-actions:permanent}{identifiant permanent}.
\end{definition}

Le \omterm{def:rs:universe-symbology-and-corporate-actions:master}{référentiel des titres} répond, à toute date, aux trois questions que pose toute table de recherche~: quel titre ce ticker désignait-il
à cette date, quel était alors le ticker de ce titre, et quels titres étaient cotés. Il se construit à partir d'événements et non à partir d'un instantané~:
une cotation crée un identifiant, un renommage met fin à une correspondance et en ouvre une autre, une radiation clôt le titre et consigne la manière dont il a pris fin,
et une division d'actions ou un dividende devient un facteur d'ajustement (livre~1, chapitre~8) rattaché à l'identifiant, avec ses dates d'annonce et
ses dates ex.

Le marché simulé montre ce que produit une jointure sur les tickers. Sur dix ans, il compte 1\,471 cotations, dont 471 prennent fin~: 113 pour
cause de performance (un \omterm{def:rs:universe-symbology-and-corporate-actions:delisting}{rendement de radiation} de $-30\%$) et 358 par fusions en numéraire (avec une prime)~; 158 tickers sont réattribués. Un fichier de cours indexé
sur les tickers contient donc 164 rendements qui courent du dernier cours d'une société au premier cours d'une autre, avec une taille absolue moyenne
de 164~\%. Le ticker KMKN en est un exemple (\cref{fig:rs:universe-symbology-and-corporate-actions:splice})~: son premier détenteur
a été négocié pour la dernière fois à 4.23 avant sa radiation~; trois mois plus tard, une nouvelle société a été cotée sous KMKN à 39.55. Dans le fichier des tickers, KMKN
a «~progressé~» de 835~\% en trois mois.

\begin{omfigure}
\begin{tikzpicture}
\begin{axis}[width=11cm,height=4.8cm,xlabel={mois},ylabel={cours},tick label style={font=\small},
  label style={font=\small},xmin=0,xmax=120,ymin=0,grid=major,grid style={gray!25},legend columns=2,
  legend style={font=\footnotesize,at={(0.5,-0.32)},anchor=north,draw=none,fill=none,
  /tikz/every even column/.append style={column sep=8pt}},table/col sep=comma,unbounded coords=jump]
\addplot[omDef,very thick] table[x=month,y=old]{figdata/research/04-universe-symbology-and-corporate-actions/splice.csv};
\addplot[omThm,very thick] table[x=month,y=new]{figdata/research/04-universe-symbology-and-corporate-actions/splice.csv};
\legend{premier détenteur de KMKN,second détenteur de KMKN}
\end{axis}
\end{tikzpicture}

\omcaption{Le ticker KMKN dans le fichier de cours indexé sur les tickers du marché simulé~: le premier détenteur tombe à 4.23 et est radié
au mois 36~; une nouvelle société est cotée sous le même ticker au mois 39 à 39.55. Une table indexée sur le ticker enregistre un rendement de
$+835\%$ entre les deux. Données~: module du chapitre, avec germe fixé.}\label{fig:rs:universe-symbology-and-corporate-actions:splice}
\end{omfigure}

\section{Ajuster sans fuite}

Les cours sont ajustés de sorte que les divisions et les dividendes n'apparaissent pas comme des rendements~: la série ajustée est le cours négocié divisé par
le produit des facteurs d'ajustement de toutes les opérations ultérieures (livre~1, chapitre~8). L'ajustement est exact pour les rendements et faux pour
les niveaux, et la différence fuit.

\begin{proposition}[Ce que les cours ajustés peuvent et ne peuvent pas servir à calculer]\label{prop:rs:universe-symbology-and-corporate-actions:adjusted}
Les rendements calculés à partir d'une série ajustée sont les rendements totaux de la détention du titre. Un niveau de la série ajustée à une
date passée, en revanche, dépend des opérations survenues \emph{après} cette date~: une règle appliquée à ce niveau à la date $t$ utilise une information
postérieure à $t$.
\end{proposition}

\begin{proof}
Entre deux opérations, l'ajustement divise tous les cours par la même constante, qui s'annule dans un rapport~; à la date ex d'une opération, le
facteur supprime la variation mécanique du cours. Le niveau en $t$ est le cours négocié divisé par le produit des facteurs des
opérations postérieures à $t$, qui ne sont pas connues en $t$.
\end{proof}

Un seuil de prix est le cas classique. Un univers de recherche exclut les actions inférieures à 5~USD, et le seuil est appliqué aux cours
ajustés d'un fichier de fournisseur. Une action qui se divise ensuite par deux avait, dans le fichier ajusté, la moitié de son cours négocié~: à 8.65, elle se négociait
nettement au-dessus du seuil, mais son cours ajusté, 4.32, se situe en dessous
(\cref{fig:rs:universe-symbology-and-corporate-actions:floor}). Pourquoi s'est-elle divisée~? Parce que son cours a monté. Le seuil appliqué aux cours
ajustés exclut donc, dès le début de l'univers, des actions sélectionnées pour leur hausse future. Dans le marché simulé, il retire à tort
0.3~\% des observations mois-titre situées au-dessus du seuil, et les titres ainsi retirés rapportent 71~\% sur les douze mois suivants, contre 15~\%
pour l'univers. Une stratégie qui paraît bonne, ou mauvaise, pour l'achat d'actions bon marché pourrait bien mesurer cette fuite.

\begin{omfigure}
\begin{tikzpicture}
\begin{semilogyaxis}[width=11cm,height=5cm,xlabel={mois},ylabel={cours (échelle log)},tick label style={font=\small},
  label style={font=\small},xmin=0,xmax=120,ymin=2,ymax=300,ytick={2,5,10,20,50,100,200},
  yticklabels={2,5,10,20,50,100,200},grid=major,grid style={gray!25},legend columns=3,
  legend style={font=\footnotesize,at={(0.5,-0.32)},anchor=north,draw=none,fill=none,
  /tikz/every even column/.append style={column sep=8pt}},table/col sep=comma]
\addplot[omDef,very thick] table[x=month,y=traded]{figdata/research/04-universe-symbology-and-corporate-actions/adjfloor.csv};
\addplot[omThm,very thick,dashed] table[x=month,y=adjusted]{figdata/research/04-universe-symbology-and-corporate-actions/adjfloor.csv};
\addplot[black!70,thick,sharp plot] coordinates {(0,5) (120,5)};
\legend{cours négocié,ajusté pour la division ultérieure,seuil de 5 USD}
\end{semilogyaxis}
\end{tikzpicture}

\omcaption{Une action simulée qui se divise par deux au mois 57. Avant la division, son cours négocié est supérieur à 5~USD, mais son
cours ajusté des divisions, qui divise par le facteur ultérieur, commence en dessous~: un seuil appliqué aux cours ajustés écarte le titre pour
avoir monté par la suite. Données~: module du chapitre, avec germe fixé.}\label{fig:rs:universe-symbology-and-corporate-actions:floor}
\end{omfigure}

La règle qui en découle est simple. Calculer les rendements à partir des cours ajustés~; appliquer chaque règle portant sur un niveau (un seuil de prix, un seuil de capitalisation, une
distance à un plus haut passé) au cours négocié ce jour-là, avec le nombre d'actions en circulation connu ce jour-là.

\section{Radiations et rendement de radiation}

\begin{definition}[Rendement de radiation]\label{def:rs:universe-symbology-and-corporate-actions:delisting}
Le \emph{rendement de radiation}\index{rendement de radiation} d'un titre est le rendement entre son dernier cours de négociation régulière et la valeur
que ses détenteurs ont effectivement reçue ou pouvaient réaliser~: le numéraire d'une fusion, le cours des actions sur le marché où elles ont continué
à se négocier, ou ce qu'une faillite a laissé.
\end{definition}

Le \omterm{def:rs:universe-symbology-and-corporate-actions:delisting}{rendement de radiation} est le dernier rendement de l'historique d'un titre, et c'est celui qui manque le plus souvent. Pour une fusion en numéraire, il est
connu et généralement faible (le cours avait convergé vers le prix de l'offre). Pour une radiation pour motif de performance, il est fort et négatif~: Shumway
(1997) a constaté que les rendements des actions de la NYSE et de l'AMEX radiées pour motif de performance manquaient en grande partie dans la base de recherche
américaine usuelle et, en reconstituant 71~\% d'entre elles à partir des cours de gré à gré, une moyenne de $-30~\%$. Un backtest qui laisse la position
disparaître à son dernier cours régulier a traité $-30\%$ comme zéro, précisément pour les positions qui perdaient. Avec un \omterm{def:rs:universe-symbology-and-corporate-actions:master}{référentiel des titres}
la règle est mécanique~: chaque radiation a un motif et un rendement~; si le rendement n'est pas connu, une estimation par motif est utilisée et
signalée, jamais zéro.

\section{Tutoriel~: un marché aux noms changeants}\label{tut:rs:universe-symbology-and-corporate-actions:master}

\begin{tutorial}
\textbf{Objectif.} Simuler dix ans de cotations, construire leur \omterm{def:rs:universe-symbology-and-corporate-actions:master}{référentiel des titres} et leurs univers, et mesurer ce que les tickers et les cours ajustés
font à une table de recherche. \textbf{État final~:} \cref{fig:rs:universe-symbology-and-corporate-actions:churn,fig:rs:universe-symbology-and-corporate-actions:splice,fig:rs:universe-symbology-and-corporate-actions:floor}~;
164 rendements raccordés~; rotation de 11.5~\% et de 3.1~\%~; les 71~\% des titres retirés contre 15~\%.
\begin{tutsteps}
\item \textbf{Résoudre à une date.} Un ticker renvoie à un seul \omterm{def:rs:universe-symbology-and-corporate-actions:permanent}{identifiant permanent} à une date~; une cotation a un ticker à une date,
ou aucun avant sa cotation et après sa radiation.
\omcode{code/firm/secmaster/firm_secmaster.py}{73}{89}{Résolution à une date, dans les deux sens.}\label{lst:rs:universe-symbology-and-corporate-actions:resolve}
\item \textbf{Univers.} Classer par valeur négociée connue chaque mois~; conserver les membres dans le rang de sortie~; compléter avec les meilleurs
non-membres.
\omcode{code/firm/secmaster/firm_secmaster.py}{107}{124}{Une règle d'appartenance à l'univers avec tampon.}\label{lst:rs:universe-symbology-and-corporate-actions:buffer}
\item \textbf{Les deux fuites.} Construire le fichier de cours indexé sur les tickers avec \texttt{ticker\_panel} et compter les rendements qui couvrent deux
cotations~; appliquer le seuil de 5~USD aux cours ajustés puis aux cours négociés avec \texttt{price\_filter\_leak}.
\end{tutsteps}
\textbf{Pour aller plus loin.} Fixer \texttt{reuse\_prob=0} et vérifier que les rendements raccordés disparaissent~; ajouter un historique minimal de
douze mois à la règle d'univers et mesurer combien de nouvelles cotations elle retarde.
\end{tutorial}

\section{Construction~: le référentiel des titres}\label{bld:rs:universe-symbology-and-corporate-actions:secmaster}

\begin{build}
\bfield{Rôle} La table de référence unique de la petite entreprise~: chaque cours, caractéristique, position et transaction des chapitres suivants renvoie à un
\omterm{def:rs:universe-symbology-and-corporate-actions:permanent}{identifiant permanent}, et chaque symbole est résolu ici à la date d'utilisation. L'univers synthétique du chapitre~5 et
les moteurs de backtest de la partie~IV l'utilisent.
\bfield{Interface} \texttt{SecurityMaster().list(pid, ticker, date)}, \texttt{rename(pid, ticker, date)},
\texttt{delist(pid, date, reason, ret)}~; \texttt{resolve(ticker, date)}, \texttt{ticker(pid, date)}, \texttt{listed(date)}, \texttt{spans(pid)},
\texttt{delisting(pid)}, \texttt{reused()}~; \texttt{buffered\_universe(dates, scores, n, exit\_rank)}~; \texttt{churn(universe)}.
\bfield{Règles} Un \omterm{def:rs:universe-symbology-and-corporate-actions:permanent}{identifiant permanent} n'est utilisé qu'une fois~; un ticker n'est détenu que par une seule cotation à une même date~; un renommage prend effet à sa
date (l'ancien ticker n'est plus détenu ce jour-là), une radiation à sa date est le dernier jour détenu~; les univers ne lisent que les scores de leur
date.
\bfield{Tests d'acceptation} \texttt{code/firm/secmaster/tests/}~: un \omterm{def:rs:universe-symbology-and-corporate-actions:permanent}{changement de ticker} résolu des deux côtés de sa date~; un ticker réutilisé
résolu vers deux cotations, avec la radiation enregistrée~; un ticker détenu refusé~; un tampon réduit la rotation.
\bfield{Pour aller plus loin} Opérations sur titres rattachées aux identifiants, avec leurs dates d'annonce (en enveloppant \texttt{firm.corpactions} du livre~1)~;
plusieurs types d'identifiants externes, chacun avec ses propres plages de dates~; des fusions qui font passer une cotation dans une autre (une fusion en actions).
\end{build}

\begin{omsources}
\item Meta Platforms, Inc., ``Meta Platforms, Inc. to change ticker symbol to `META' on June 9'', communiqué de presse, 31 mai 2022 (dépôt auprès de la SEC,
pièce 99.1).
\item General Motors Company, déclaration d'enregistrement Form S-1/A, novembre 2010 (cotation sous ``GM''~; actions de
Motors Liquidation Company négociées sous MTLQQ).
\item ISO TC 68, ``What is ISIN?''~; CUSIP Global Services, descriptions des identifiants~; OpenFIGI, ``About FIGI''.
\item T.~Shumway, ``The delisting bias in CRSP data'', \emph{Journal of Finance} 52(1), 1997.
\end{omsources}

\section{Exercices}

\begin{exercise}[$\star$]\label{exo:rs:universe-symbology-and-corporate-actions:1}
Une action se négocie à 12~USD et fait l'objet d'une division par trois un an plus tard. Quel est son cours ajusté aujourd'hui pour cette date, et un
plancher de 5~USD sur les cours ajustés l'inclut-il~?
\end{exercise}

\begin{exercise}[$\star$]\label{exo:rs:universe-symbology-and-corporate-actions:2}
Une table indexée sur les tickers contient une société radiée à 0.80~USD et, deux mois plus tard, une autre cotée sous le même ticker
à 24~USD. Quel rendement la table affiche-t-elle~?
\end{exercise}

\begin{exercise}[$\star$]\label{exo:rs:universe-symbology-and-corporate-actions:3}
Un univers équipondéré de 500 noms remplace 11.5~\% de ses membres par mois~; chaque nom remplacé est vendu et son remplaçant
acheté, à 20 points de base de la valeur échangée de chaque côté. Quel est le coût annuel de la seule rotation des appartenances, en part du capital~?
Et avec le tampon à 3.1~\%~?
\end{exercise}

\begin{exercise}[$\star\star$]\label{exo:rs:universe-symbology-and-corporate-actions:4}
Sur un an, 2~\% des noms d'un univers sont radiés pour mauvaise performance, sans \omterm{def:rs:universe-symbology-and-corporate-actions:delisting}{rendement de radiation} dans les données, et une stratégie les détient
à poids moyen. De combien l'omission d'un \omterm{def:rs:universe-symbology-and-corporate-actions:delisting}{rendement de radiation} de $-30\%$ surestime-t-elle le rendement équipondéré~?
\end{exercise}

\begin{exercise}[$\star\star$]\label{exo:rs:universe-symbology-and-corporate-actions:5}
Pourquoi un rang de sortie plus large que la taille cible réduit-il davantage la rotation qu'il ne réduit la liquidité~? Donner un argument en termes de la façon dont les rangs
évoluent d'un mois à l'autre.
\end{exercise}

\begin{exercise}[$\star\star$]\label{exo:rs:universe-symbology-and-corporate-actions:6}
Un signal utilise la distance du cours à son plus haut sur 52~semaines. Sur quelle série doit-il être calculé~: cours échangés, cours ajustés des divisions ou
cours ajustés du rendement total~? Que se passe-t-il avec chacune des autres~?
\end{exercise}

\begin{exercise}[$\star\star\star$]\label{exo:rs:universe-symbology-and-corporate-actions:7}
\emph{Programmation.} Relancer le marché simulé avec \texttt{reuse\_prob} égal à 0, 0.35 et 0.7. Comment évoluent le nombre de rendements raccordés et
leur taille absolue moyenne et médiane~? Pourquoi la moyenne est-elle nettement supérieure à la médiane~?
\end{exercise}

\begin{exercise}[$\star\star\star$]\label{exo:rs:universe-symbology-and-corporate-actions:8}
\emph{Trouver la faille.} «~Notre univers est constitué des actions figurant dans le fichier courant de notre fournisseur, dont le cours ajusté dépasse 5~USD et dont l'historique
couvre au moins trois ans. Nous le reconstruisons chaque mois depuis 2010 à partir du même fichier.~»
\end{exercise}

\section{Problème~: le ticker qui a changé de mains}

\begin{problem}[{Problème du week-end --- ce que font une table indexée sur les tickers et un plancher de cours ajustés à un univers}]\label{pb:rs:universe-symbology-and-corporate-actions:1}
Le marché simulé du chapitre~: 1\,000 cotations à tout instant sur 120 mois, radiations pour mauvaise performance à un rendement logarithmique cumulé
de $-1.6$ avec $-30~\%$, fusions en numéraire à 0.3~\% par mois avec une prime de 25~\%, divisions au-delà de 200~USD, renommages à 0.2~\% par mois, et
tickers libérés réutilisés par 35~\% des nouvelles cotations, graine~4.

\medskip\noindent\textbf{Partie I --- Le marché.}
\begin{enumerate}
\item Combien y a-t-il de cotations sur les dix années, et combien prennent fin~?
\item Combien prennent fin pour mauvaise performance, et combien par fusion~?
\item Combien y a-t-il de divisions, et sur combien de noms~?
\item Combien de tickers sont réutilisés, sur combien de tickers ayant jamais servi~?
\end{enumerate}

\medskip\noindent\textbf{Partie II --- Le fichier indexé sur les tickers.}
\begin{enumerate}[resume]
\item Combien de rendements du fichier par ticker couvrent deux cotations~?
\item Quelle est leur taille absolue moyenne~?
\item Quel rendement le fichier affiche-t-il pour KMKN entre les mois 36 et 39~?
\item Quelle part des couples (mois, ticker) de l'univers tamponné ont un historique de douze mois mêlant deux cotations~?
\item Pourquoi cette part est-elle très inférieure à celle des tickers réutilisés~?
\end{enumerate}

\medskip\noindent\textbf{Partie III --- Le plancher sur cours ajustés.}
\begin{enumerate}[resume]
\item Quelle part des observations mois-nom au-dessus de 5~USD un plancher sur cours ajustés écarte-t-il à tort~?
\item Quel est le rendement sur les douze mois suivants des noms écartés, et celui de l'univers~?
\item Pourquoi les noms écartés sont-ils des gagnants futurs~?
\item Quels sont les cours échangé et ajusté de l'action exemple au mois~0, et quand se produit sa division~?
\item Quelles autres règles fondées sur un niveau partagent ce défaut~?
\end{enumerate}

\medskip\noindent\textbf{Partie IV --- Univers.}
\begin{enumerate}[resume]
\item Quelles sont les rotations mensuelles de l'univers de 500 noms sans rang de sortie, puis avec un rang de sortie de 600~?
\item Combien coûte le tampon en liquidité~?
\item D'où proviennent les derniers rendements des noms radiés dans la simulation, et dans les données réelles~?
\item Énoncer le \emph{résultat nommé}~: le nombre et la taille absolue moyenne des rendements raccordés, et le rendement à terme des noms écartés par le
plancher sur cours ajustés, comparé à celui de l'univers.
\item Quelles deux tables un \omterm{def:rs:universe-symbology-and-corporate-actions:master}{référentiel des titres} rend-il inutiles~?
\item En une phrase~: pourquoi une table de recherche ne doit-elle jamais être indexée sur un ticker~?
\end{enumerate}
\end{problem}

\section{Questions d'entretien}

\begin{interviewq}[$\star$ \iqroles{développeur, chercheur}]\label{iq:rs:universe-symbology-and-corporate-actions:1}
Pourquoi un ticker n'est-il pas une bonne clé pour une table de recherche~? Qu'utiliseriez-vous à la place~?
\end{interviewq}

\begin{interviewq}[$\star\star$ \iqroles{chercheur}]\label{iq:rs:universe-symbology-and-corporate-actions:2}
Vous filtrez votre univers avec un cours minimal de 5~USD sur des données ajustées. Quel biais venez-vous d'introduire~?
\end{interviewq}

\begin{interviewq}[$\star\star$ \iqroles{développeur}]\label{iq:rs:universe-symbology-and-corporate-actions:3}
Concevez le schéma d'un \omterm{def:rs:universe-symbology-and-corporate-actions:master}{référentiel des titres} qui répond efficacement à la question «~quel titre portait le ticker X à la date D~?~».
\end{interviewq}

\begin{interviewq}[$\star\star$ \iqroles{chercheur, trader}]\label{iq:rs:universe-symbology-and-corporate-actions:4}
Comment choisissez-vous la taille et la règle de rééquilibrage d'un univers d'actions pour une stratégie quotidienne~?
\end{interviewq}

\begin{interviewq}[$\star\star$ \iqroles{chercheur, risque}]\label{iq:rs:universe-symbology-and-corporate-actions:5}
Votre backtest détient une action radiée après le dépôt d'une procédure de faillite. Quel rendement comptabilisez-vous, et pourquoi~?
\end{interviewq}

\begin{interviewq}[$\star\star\star$ \iqroles{chercheur}]\label{iq:rs:universe-symbology-and-corporate-actions:6}
Démontrez que les rendements calculés à partir de cours ajustés des divisions égalent les rendements de détention, et exhiber une règle sur les niveaux ajustés qui regarde vers le
futur.
\end{interviewq}
